% =========================================================================
% Chapter 1 — State of the Art
% =========================================================================
\chapter{State of the Art}

\chapintro
The state of the art is the explicit knowledge about software engineering and applied machine
learning that is documented from various sources. Building an AI-powered platform for
short-form video requires bringing together domains as different as multimodal deep learning,
classical image processing, full-stack web engineering and OAuth integration with three
different social-network APIs, and a mature understanding of what creators actually do with
their tools. This chapter starts by presenting the project context, then surveys existing
solutions on the market, derives the functional and non-functional requirements, identifies
the actors of the system, and finally specifies the working methodology adopted for the
project.

% -------------------------------------------------------------------------
\section{Project context and host environment}
\label{sec:context}

This end-of-studies project was carried out within the framework of an applied machine
learning and full-stack engineering internship. The host structure focuses on building
production-grade SaaS tools for content creators and small marketing teams: it operates a
TypeScript/Node.js back-end stack, a React front-end stack and a Python machine learning
sub-stack, all of which were used during this work. The project was scoped to deliver a
complete, working MVP of an AI-powered short-form video platform within roughly four months,
with three of those months dedicated to research, dataset construction and modelling, and the
last month dedicated to integration and the dashboard.

\begin{figure}[H]
  \centering
  \placeholder[10cm]{Logo of the host structure}
  \caption{Host structure logo}
  \label{fig:host_logo}
\end{figure}

% -------------------------------------------------------------------------
\section{Preliminary study}
\label{sec:prelim}

Preliminary system study is the first stage of the system development life cycle in which we
study the viability of the project. We start by a study of the existing where we analyse the
solutions already available on the market, list their drawbacks, then suggest a more suitable
solution adapted to the needs of short-form video creators.

\subsection{Study of existing solutions}
\label{sec:existing}

To ensure the quality of the final system, conducting a study of existing tools is crucial.
This phase allows us to analyse the existing solutions, their pros and cons, and make sure we
avoid those drawbacks in our own product.

\subsubsection{Analysis of existing solutions}
\label{sec:existing_analysis}

In this part we take a thorough look at the various existing solutions in the short-form
video market, focusing on their advantages and disadvantages. The four solutions we conduct
our study upon are \textbf{VidIQ}, \textbf{TubeBuddy}, \textbf{CapCut AI} and \textbf{Descript}.

\paragraph{VidIQ} is a YouTube-centric SaaS that analyses the performance of \emph{already
published} videos. It provides keyword research, tag suggestions, thumbnail analysis,
competitor monitoring and a daily ``score" derived from heuristics over the YouTube Data API.
Its strengths are a clean dashboard and an extensive keyword database; its weakness, from the
point of view of a Shorts creator, is that all signals are post-publication and largely
text-only --- the platform does not analyse audio or visual frames.

\begin{figure}[H]
  \centering
  \placeholder[12cm]{Screenshot of the VidIQ dashboard}
  \caption{Existing solution \textquotedblleft VidIQ\textquotedblright}
  \label{fig:vidiq}
\end{figure}

\paragraph{TubeBuddy} is a browser extension and SaaS that overlays YouTube Studio with
A/B-tested thumbnails, keyword explorer and a ``Best Time to Publish" estimator. Like VidIQ
it is built on top of YouTube's published-video API and offers no insight into the multimodal
content of the clip itself, no dynamic captioning and no enhancement.

\begin{figure}[H]
  \centering
  \placeholder[12cm]{Screenshot of the TubeBuddy interface}
  \caption{Existing solution \textquotedblleft TubeBuddy\textquotedblright}
  \label{fig:tubebuddy}
\end{figure}

\paragraph{CapCut AI} is a desktop and web video editor with AI-powered automatic captions,
auto-cut and a few one-click filters. It is the strongest of the four on the
\emph{post-production} side: captions are reasonable, the export is fast and the filters cover
basic colour grading. However, it is purely an editor: it neither predicts performance, nor
runs super-resolution or face restoration, nor schedules publishing across platforms.

\begin{figure}[H]
  \centering
  \placeholder[12cm]{Screenshot of the CapCut AI editor}
  \caption{Existing solution \textquotedblleft CapCut AI\textquotedblright}
  \label{fig:capcut}
\end{figure}

\paragraph{Descript} is a transcription-driven editor that turns a video into a Word-like
document and lets the user edit the video by editing the text. The captioning experience is
strong; the application has a basic ``Studio Sound" voice cleaner but no super-resolution, no
face restoration and no virality model.

\begin{figure}[H]
  \centering
  \placeholder[12cm]{Screenshot of the Descript editor}
  \caption{Existing solution \textquotedblleft Descript\textquotedblright}
  \label{fig:descript}
\end{figure}

\subsubsection{Critique of existing solutions}
\label{sec:existing_critique}

The four solutions surveyed have the following advantages and limitations:

\begin{itemize}
  \item \textbf{VidIQ and TubeBuddy} are equipped with rich keyword and competitor analytics
        and integrate cleanly with YouTube's APIs. Both, however, work \emph{only after} a
        video is published, are limited to text-based signals and ignore the visual and
        audio content entirely. They cannot tell a creator that the visual hook in the first
        second is too low-contrast or that the audio energy is below the threshold associated
        with viral clips.
  \item \textbf{CapCut AI} delivers strong basic post-production (cut, captions, filters) and
        a pleasant user experience, but it does not predict performance, does not perform
        super-resolution beyond simple upscale, has no face-aware restoration, and does not
        explain its filter choices. The captioning algorithm uses a fixed font size, which
        causes long words to overflow on portrait-mode 1080$\times$1920 video.
  \item \textbf{Descript} is excellent for text-driven editing but is fundamentally a
        post-production tool; it does not address the visual-quality gap and provides no
        predictive feedback before publication.
\end{itemize}

\noindent\textbf{Bottom line.} None of these tools offers --- in a single integrated platform
--- multimodal pre-publication predictive scoring, font-aware automatic captioning that scales
per caption based on the longest word, super-resolution and face restoration with quality
metrics, and a multi-platform publishing layer. That gap is what \textbf{BViral Control
Center} is designed to fill.

\subsection{Suggested solution}
\label{sec:suggested}

We suggest a unified, web-based platform built around four AI pillars and a single dashboard
that pipes them together end-to-end:

\begin{itemize}
  \item \textbf{Pillar I --- Pre-Posting Predictive Agent (AI Coach)}: a multimodal pipeline
        that ingests a raw video, extracts visual embeddings with CLIP, audio descriptors
        with Librosa, transcript embeddings with Whisper + MiniLM and Llama-3 generated
        semantic features (tone, emotion, hook quality), reduces the resulting concatenation
        to a 473-dimensional vector via PCA, predicts an expected performance score with an
        XGBoost regressor, and returns SHAP explanations so the creator knows \emph{why} a
        clip is or isn't predicted to perform.
  \item \textbf{Pillar II --- Dynamic Auto-Captioning Engine}: a faster-whisper-based
        transcription pipeline that produces word-level timestamps, regroups them into
        2-word ``flashes", and renders them as ASS subtitles where each line gets a
        per-caption font size computed by a safe-zone algorithm (described in
        Section~\ref{sec:pillar2}). FFmpeg burns the captions in with \texttt{libx264} at
        CRF~18 and \texttt{+faststart} for instant web playback.
  \item \textbf{Pillar III --- AI Video Enhancement Pipeline}: a 12-stage pipeline that
        assesses video quality with OpenCV (blur, noise, low-light, face presence), runs
        classical filters (median, bilateral, unsharp mask), executes Real-ESRGAN
        $2\times$/$4\times$ super-resolution, applies GFPGAN face restoration when faces are
        detected, reassembles with audio, and computes PSNR / SSIM / LPIPS / flicker scores
        for quality assurance.
  \item \textbf{Pillar IV --- AI Studio (Web Dashboard)}: a React/Vite + Fastify + Drizzle +
        PostgreSQL + Supabase + BullMQ application that orchestrates the three pillars,
        manages YouTube / TikTok / Meta OAuth (with PKCE for TikTok), schedules cross-platform
        publishing, and surfaces post-publication analytics.
\end{itemize}

The platform should provide an attractive, ergonomic dashboard that makes the underlying
machine learning approachable, and a clear set of actor-facing flows for the two main types
of users introduced in Section~\ref{sec:actors}.

% -------------------------------------------------------------------------
\section{Requirement specification}
\label{sec:requirements}

A software requirements specification is the description of the system to be developed and
its functional and non-functional requirements. It forms the basis for the agreement between
stakeholders on the way the software product should function.

\subsection{Identifying major actors in the system}
\label{sec:actors}

The system has two main actors:

\begin{itemize}
  \item \textbf{Content Creator.} The primary user. Uploads short-form videos, consumes the
        AI Coach's score and SHAP advice, runs the captioning and enhancement pipelines,
        connects social accounts and schedules publication. The creator owns their videos
        through PostgreSQL Row Level Security policies.
  \item \textbf{Administrator.} A privileged user that can supervise the multi-tenant
        platform: manage users, monitor processing-queue health, review alerts (errors,
        copyright strikes, OAuth token expiry), inspect aggregate analytics and re-run failed
        jobs. Administrators have the \texttt{role = `admin'} flag in the \texttt{users}
        table.
\end{itemize}

\subsection{Functional requirements}
\label{sec:func_req}

Our solution should satisfy the following functional requirements.

\noindent\textbf{The Content Creator:}
\begin{itemize}
  \item Authenticate against the platform using a Supabase-issued JWT and view a personalised
        dashboard.
  \item Upload an MP4 or MOV video to Supabase Storage and follow the upload progress.
  \item Trigger the Pre-Posting Predictive Agent on a video and obtain a virality score plus
        SHAP-based feature explanations.
  \item Trigger the Dynamic Auto-Captioning Engine on a video and download the captioned
        version, the SRT file or the styled ASS file.
  \item Trigger the AI Video Enhancement Pipeline with one of three modes
        (\texttt{fast}, \texttt{hd}, \texttt{ultra}) and receive a side-by-side preview plus
        the final enhanced video.
  \item Connect a YouTube, TikTok or Meta account through OAuth.
  \item Schedule a post for one or several connected accounts with a future publication time.
  \item Browse post-publication analytics for owned posts (views, likes, comments, shares,
        revenue).
\end{itemize}

\noindent\textbf{The Administrator:}
\begin{itemize}
  \item Authenticate with elevated permissions and view the global health dashboard.
  \item Manage user accounts (create, suspend, change role).
  \item Resolve alerts (errors, copyright issues, expired refresh tokens).
  \item Inspect cross-tenant aggregate analytics and queue lengths.
\end{itemize}

\subsection{Non-functional requirements}
\label{sec:nonfunc_req}

Non-functional requirements define system attributes that are mostly hidden from the end
user. The platform should obey the following rules:

\begin{itemize}
  \item \textbf{Performance.} GPU-side AI inference must complete in under 60 seconds per
        minute of input video on an NVIDIA T4 / L4-class GPU for the captioning and HD
        enhancement modes. The Fastify API must respond to dashboard CRUD calls in under
        200~ms p95.
  \item \textbf{Security.} Authentication uses a JWT bearer token issued by Supabase. OAuth
        refresh tokens stored in the \texttt{accounts} table are encrypted at rest with
        AES-256-GCM using a 32-byte key supplied via the \texttt{TOKEN\_ENCRYPTION\_KEY}
        environment variable; the API server refuses to boot without it. PostgreSQL Row
        Level Security policies ensure that a user can only read or modify their own data.
  \item \textbf{Scalability.} The publishing pipeline runs as a BullMQ-on-Redis queue with a
        dedicated worker process so that bursty social-network publishing does not back up
        the API event loop. The same model is used for video processing and analytics
        refresh.
  \item \textbf{Portability.} The dashboard, the API and the workers run on Node.js~24 and
        ship as standalone bundles; the Python AI services run inside Docker containers and
        can be deployed on any GPU-enabled host (CUDA optional --- the enhancement pipeline
        also has a parallel-CPU fallback).
  \item \textbf{Maintainability.} The codebase is fully typed (TypeScript end-to-end, Zod
        schemas shared between client and server, Python type hints) and follows a
        monorepo / pnpm-workspace layout that keeps domain-specific dependencies isolated.
  \item \textbf{Ergonomics.} The dashboard uses Tailwind CSS and shadcn/ui components for a
        modern, dark-themed interface with smooth transitions, drag-and-drop video upload
        and real-time toast notifications.
  \item \textbf{Interoperability.} A vendored Next.js video editor (BVIRAL Video Editor) is
        embedded as an iframe so that creators can perform manual timeline edits without
        leaving the platform.
\end{itemize}

% -------------------------------------------------------------------------
\section{Working methodology}
\label{sec:methodology}

When developing this solution we had to take into account all the necessary constraints of a
modern AI product: model experimentation, dataset evolution, dashboard ergonomics, deployment
and continuous integration. We chose an \textbf{Agile / Scrum} methodology with two-week
sprints rather than the V-Model used in many pure-engineering projects, because machine
learning work is intrinsically iterative: each modelling sprint generates new feature ideas,
and the dashboard requirements evolve as soon as the predictive scores become trustworthy
enough to surface to the user.

Each sprint had four ceremonies: a sprint planning at the beginning, a daily 15-minute
stand-up, a sprint review at the end, and a short retrospective. The planning produced a
sprint backlog made of two kinds of items: \emph{engineering tickets} (e.g. ``add the
\texttt{accounts} TikTok connect route") and \emph{ML experiments} (e.g. ``train an XGBoost
regressor on top of the 473-dim PCA features and report R\textsuperscript{2} on the
held-out test split"). Engineering tickets followed a definition of done that included
unit-tested code, type-check pass and a Pull Request reviewed by the supervisor; ML
experiments followed a definition of done that included a notebook with the metrics, a
saved model artefact and a one-line summary in the experiments tracker.

\begin{figure}[H]
  \centering
  \placeholder[12cm]{Diagram of the Scrum sprint cycle}
  \caption{Agile / Scrum life cycle adopted for the project}
  \label{fig:scrum}
\end{figure}

% -------------------------------------------------------------------------
\section{Gantt diagram}
\label{sec:gantt}

This work was carried out over a period of four months. The distribution of tasks across the
project timeline is illustrated by the Gantt chart below.

\begin{figure}[H]
\centering
\begin{ganttchart}[
    hgrid,
    vgrid,
    x unit=0.55cm,
    y unit chart=0.55cm,
    bar/.append style={fill=blue!40},
    bar height=0.55,
    title height=0.7,
    title label font=\footnotesize\bfseries,
    bar label font=\footnotesize,
    group label font=\footnotesize\bfseries,
    milestone label font=\footnotesize\itshape,
  ]{1}{16}
  \gantttitle{Month 1}{4}
  \gantttitle{Month 2}{4}
  \gantttitle{Month 3}{4}
  \gantttitle{Month 4}{4}\\
  \gantttitle{Week}{16}\\
  \ganttbar{Specification \& state of the art}{1}{2}\\
  \ganttbar{Dataset collection (TikTok scraping)}{2}{4}\\
  \ganttbar{Pillar~I --- multimodal feature pipeline}{3}{6}\\
  \ganttbar{Pillar~I --- XGBoost training \& SHAP}{5}{7}\\
  \ganttbar{Pillar~II --- captioning engine}{6}{8}\\
  \ganttbar{Pillar~III --- enhancement pipeline}{7}{10}\\
  \ganttbar{Pillar~IV --- AI Studio dashboard}{9}{14}\\
  \ganttbar{Fastify API \& Drizzle schema}{8}{13}\\
  \ganttbar{OAuth integrations \& BullMQ workers}{11}{14}\\
  \ganttbar{Integration testing}{13}{15}\\
  \ganttbar{Final report writing}{14}{16}\\
  \ganttmilestone{MVP demo}{16}
\end{ganttchart}
\caption{Gantt diagram of the project}
\label{fig:gantt}
\end{figure}

% -------------------------------------------------------------------------
\chapconclusion
In this chapter we performed the study of the existing solutions in the short-form video
optimisation market, identified the four-pillar gap that none of them addresses, suggested an
integrated solution, listed the major actors and specified both the functional and the
non-functional requirements. We then justified an Agile / Scrum methodology over the more
classical V-Model and presented the Gantt diagram. The next chapter will dive into the
detailed conception of the solution through UML diagrams and the specification of the AI
pipelines and dataset.

\cleardoublepage
